% !PRINTED 45
\weberpage{45}
%<<CONT

3. \gesperrt{The values of a character $\chi(A)$ are roots of unity whose degree is a divisor of the degree of $A$.}

After this we can easily determine all the characters. For if we set
\weq{9}{
\chi(A_1) = \omega_1,\ \chi(A_2) = \omega_2 \ldots \chi(A_\nu) = \omega_\nu,
}
then
\weq{10}{\omega_1^{a_1} = 1,\ \omega_2^{a_2} = 1, \ldots \omega_\nu^{a_\nu} = 1,}
and by (2) and (4)
\weq{11}{\chi(A) = \omega_1^{\alpha_1}\omega_2^{\alpha_2}\ldots\omega_\nu^{\alpha_\nu}.}

Conversely, if $\omega_1, \omega_2 \ldots \omega_\nu$ denotes any solution of the equations (10), then a function of $A$ is determined by (11) which, by 1., satisfies the condition (4) and is therefore a character of the group.

Each of the equations (10) has $a_1, a_2 \ldots a_\nu$ roots, and if we combine each with each, there result
\[
a_1 a_2 \ldots a_\nu = n
\]
combinations. All these combinations lead to distinct characters $\chi(A)$. For if we denote by $\omega_1', \omega_2' \ldots \omega_\nu'$ a second combination of roots of the equation (10), and if for all elements $A$
\[
\omega_1^{\alpha_1}\omega_2^{\alpha_2}\ldots\omega_\nu^{\alpha_\nu} = \omega_1'^{\alpha_1}\omega_2'^{\alpha_2}\ldots\omega_\nu'^{\alpha_\nu},
\]
then, if one sets $\alpha_1 = 1, \alpha_2 = 0 \ldots \alpha_\nu = 0$, it follows that $\omega_1 = \omega_1'$, and likewise $\omega_2 = \omega_2' \ldots \omega_\nu = \omega_\nu'$. From this follows the theorem:

4. \gesperrt{There are $n$ and no more distinct characters of an Abelian group of degree $n$, all of which are represented by the formula (11).}

If we understand by $\varepsilon_1, \varepsilon_2 \ldots \varepsilon_\nu$ a system of primitive roots of the equations (10), we can set
\weq{12}{\omega_1 = \varepsilon_1^{\beta_1},\ \omega_2 = \varepsilon_2^{\beta_2}, \ldots \omega_\nu = \varepsilon_\nu^{\beta_\nu}}
and can take the $\beta_1, \beta_2 \ldots \beta_\nu$, just like the $\alpha_1, \alpha_2 \ldots \alpha_\nu$, each from a complete system of residues with respect to the moduli $a_1, a_2 \ldots a_\nu$. Then, with $\varepsilon_1, \varepsilon_2 \ldots \varepsilon_\nu$ held fixed, we obtain all the $n$ group characters in the form
\weq{13}{\chi(A) = \varepsilon_1^{\alpha_1\beta_1}\varepsilon_2^{\alpha_2\beta_2}\ldots\varepsilon_\nu^{\alpha_\nu\beta_\nu}.}
